\subsubsection{從米勒指數找到法向}

晶面法向由倒格矢確定。採用第\ref{ch:theory}章的倒格約定，表面對應的方向為
\begin{equation}
  \mathbf G_{hkl}=h\mathbf b_1+k\mathbf b_2+l\mathbf b_3,
  \qquad
  \hat{\mathbf n}=\frac{\mathbf G_{hkl}}{\|\mathbf G_{hkl}\|}.
  \label{eq:surface-miller-normal}
\end{equation}
式中，$\mathbf b_1,\mathbf b_2,\mathbf b_3$為倒格矢，單位為$\text{\AA}^{-1}$；$h,k,l$是無量綱整數，$\hat{\mathbf n}$是單位法向量。立方晶胞的三條邊等長且正交，$(111)$法向與$\mathbf a+\mathbf b+\mathbf c$平行；本例的晶格軸沿全域性座標軸排列，因此也與全域性$(1,1,1)$平行。斜晶胞一般沒有這種對應關係，不能把米勒指數直接當作笛卡爾方向輸入吸附物的定向引數。

構建後的表面晶胞以$\mathbf a,\mathbf b$為面內向量，可用
\begin{equation}
  \hat{\mathbf n}_{\mathrm{slab}}=
  \frac{\mathbf a\times\mathbf b}{\|\mathbf a\times\mathbf b\|},
  \qquad H=|\mathbf c\cdot\hat{\mathbf n}_{\mathrm{slab}}|
  \label{eq:surface-normal-height}
\end{equation}
檢查實際法向與晶胞高度。$H$是相鄰週期表面之間沿法向的重複長度，並不總等於$\|\mathbf c\|$。把同一薄片的原子放在連續的週期映像中後，將原子位置投影到該法向，便可讀出最外兩層之間的跨度。它與構建視窗的最小厚度、完整重複單元厚度也應分別記錄。

\begin{table}[H]
\centering
\caption{表面構建中幾種長度的含義。}
\begin{tabularx}{\linewidth}{lX}
\toprule
量 & 怎樣理解與核對\\
\midrule
最小薄片厚度 & 構建時用於選擇完整重複單元的條件；不是直接指定層數。\\
原子層跨度 & 最外兩層的法向距離；檢查時先選取同一連續薄片。\\
法向胞高$H$ & 由式\eqref{eq:surface-normal-height}得到，包含薄片和週期間隔。\\
新增真空$\Delta$ & \ui{新增真空層}中沿第三晶格向量增加的長度，不是最終總真空。\\
\bottomrule
\end{tabularx}
\end{table}

對於非正交輸入，還要檢查實際輸出的法向、層數和週期間隔。斜晶胞的厚度換算與候選約簡可能受到取向處理的影響，不宜只依據輸入的最小厚度判斷結果。本章的Pt例子從立方晶胞開始，並檢查生成的表面。
